\section{Related Work}
\label{sec:related}

% NOTE: packages required: booktabs. Sources are listed in refs_related.bib; items read only through related_work_report.md are marked in that file.

\subsection{Hallucination in vision-language models}
\label{subsec:rw_vlm}

Hallucination in image-conditioned text generation was first quantified for image captioning. The CHAIR metric measures how often a captioning model mentions objects that are absent from the image, and the accompanying analysis showed that the models with the best scores on standard sentence-level metrics do not always hallucinate least \citep{rohrbach2018chair}. For large vision-language models, POPE recasts the problem as a set of polling questions on the presence of objects. Its authors report that most of the models they tested frequently describe objects that do not correspond to the image, and that objects which occur often in instruction data or together with other objects in an image are particularly prone to being hallucinated \citep{li2023pope}. Surveys of hallucination in natural language generation and in multimodal language models document the causes, the evaluation methods and the mitigation approaches proposed so far \citep{ji2023survey,bai2024hallucination}. These benchmarks treat each error as a unit of equal weight and are built on natural images, in which the presence of an object can be verified from common knowledge. Neither condition holds in medical imaging, where the verification of a statement requires a ground truth derived from expert interpretation and where the consequence of an error depends on its type.

\subsection{Claim-level and automated evaluation}
\label{subsec:rw_claims}

Evaluating a long free-text response requires a unit of analysis smaller than the response. FActScore decomposes the output of a language model into short, self-contained claims and reports the share of claims supported by a trusted reference, together with an automated estimator that closely approximates human judgements \citep{min2023factscore}. The claim is also the unit of analysis in the present work. The reference, however, is a structured image-level ground truth and not a text corpus, and the score has to account for statements that a response fails to make, which a precision-oriented measure does not capture.

Automated judgement by a language model is the usual means of scaling such evaluation. In a general-purpose setting, a strong proprietary judge was found to agree with controlled expert preferences and with crowdsourced human preferences at a rate above 80\%, which matches the agreement observed between human raters \citep{zheng2023judge}. That result concerns a large judge and general conversational preferences. It does not transfer automatically to a smaller locally hosted judge applied to medical claims, and the labelling pipeline used here (Section~\ref{sec:benchmark}) is therefore examined separately in Section~\ref{sec:validation}.

\subsection{Medical vision-language models}
\label{subsec:rw_medvlm}

Domain-adapted VLMs start from a general architecture and are further trained on biomedical image--text data. LLaVA-Med was trained to act as an assistant for biomedical images \citep{li2023llavamed}. CheXagent was developed as a foundation model for the interpretation of chest radiographs \citep{chen2024chexagent}. HuatuoGPT-Vision was built by injecting medical visual knowledge into a multimodal language model at scale \citep{chen2024huatuogptvision}. The three models differ in scope. CheXagent targets one modality, whereas the other two address medical images in general. This difference is relevant to the interpretation of the specialisation results in Section~\ref{sec:results}, because the performance of CheXagent on dermatology, histopathology and ophthalmology images measures in part the transfer of a chest radiograph model to other domains.

\subsection{Hallucination benchmarks for medical vision-language models}
\label{subsec:rw_bench}

Several benchmarks examine hallucination in medical VLMs, and they differ in task, in scope and in the way an error is graded. Table~\ref{tab:related} summarises their main properties.

\begin{table}[t]
\centering
\small
\caption{Properties of benchmarks for hallucination in medical vision-language models, as stated in the cited sources. A dash indicates that the property was not established from the sources consulted.}
\label{tab:related}
\begin{tabular}{@{}p{2.3cm}p{4.9cm}p{3.5cm}p{3.1cm}@{}}
\toprule
Work & Scope and task & Models evaluated & Severity or risk grading \\
\midrule
\citet{wu2024halluvqa} & PMC-VQA, PathVQA and VQA-RAD recast as multiple-choice questions (FAKE, NOTA, image swap) & LLaVA variants including LLaVA-Med; GPT-4-turbo-vision & --- \\
\citet{chen2024medhallmark} & Medical VQA and report generation & --- & Five levels (MediHall Score) \\
\citet{gu2024medvh} & Five tasks, including long-form response generation & --- & --- \\
\citet{yan2025medhalltune} & Over 100,000 images and over 1,000,000 instruction pairs & Medical and general-purpose VLMs & Risk level, one of four criteria \\
\citet{chang2025medheval} & Close-ended and open-ended medical VQA, radiology; three causes of hallucination & 11 VLMs; 7 mitigation methods & --- \\
\citet{wang2026hallucxr} & 856 frontal chest radiographs; three query types; eight-category taxonomy & 6 models & Three levels, reported as the share of outputs with a level-3 hallucination \\
\citet{yang2026clinhallu} & Stage-wise decomposition of reasoning & --- & --- \\
\citet{crest2025medrxiv} & 685 chest radiograph cases; report impressions & 3 VLMs & Three clinical risk levels \\
\citet{xia2024cares} & 41,000 question--answer pairs, 16 modalities; trustworthiness & --- & --- \\
This work & \Nimages{} images, 4 modalities, up to 4 query types & \Nmodels{} VLMs & Seven classes with weights \\
\bottomrule
\end{tabular}
\end{table}

Wu et al.\ adapted three public medical VQA datasets into multiple-choice questions under three scenarios: nonsensical questions, questions whose correct option is replaced by ``none of the above'', and questions paired with an unrelated image. They found that domain fine-tuning did not guarantee resistance to hallucination, since LLaVA-Med performed worse than the general model from which it was derived \citep{wu2024halluvqa}. MedVH comprises five tasks that cover the comprehension of text and visual input as well as long-form response generation, and reports that the medical models it evaluated often hallucinate more than general-purpose ones \citep{gu2024medvh}. MedHallTune provides over 100,000 images and over 1,000,000 instruction pairs, with and without hallucinated content, and evaluates models on clinical accuracy, relevance, level of detail and risk level. Its authors report that fine-tuning on the dataset reduced hallucination \citep{yan2025medhalltune}. MedHEval groups hallucinations by cause into visual misinterpretation, knowledge deficiency and context misalignment, and evaluates 11 models together with seven mitigation techniques on medical VQA data \citep{chang2025medheval}.

Work on radiology is more detailed. HalluCXR evaluates six models on 856 frontal chest radiographs from MIMIC-CXR, with CheXpert labels as ground truth. Its taxonomy has eight categories: fabrication, omission, spatial error, severity distortion, reasoning error, template overfitting, self-contradiction and detail fabrication. Each hallucination receives a severity of 1 to 3, and severity is reported as the share of outputs that contain at least one hallucination of level 3. No aggregate severity-weighted score is computed. The study reports that normal radiographs drew the most severe hallucinations \citep{wang2026hallucxr}. Five of the seven classes of the present taxonomy correspond to categories of HalluCXR, and the taxonomy should be read as an adaptation and extension of that scheme and not as an independent construction. ClinHallu decomposes the reasoning of a medical model into stages in order to locate the stage at which a hallucination arises \citep{yang2026clinhallu}. CARES assesses several dimensions of trustworthiness over 41,000 question--answer pairs from 16 modalities, and is therefore broad in modality but not specific to hallucination \citep{xia2024cares}. The finding of HalluCXR that normal images are the most hazardous is consistent with the design of the present benchmark, in which a false-premise query is applied to normal images in every modality.
% NOTE: the HalluCXR taxonomy and severity definitions were read from the full text (Table 3, Section 3.4) in this session.

\subsection{Evaluation outside radiology}
\label{subsec:rw_nonrad}

Studies of other imaging domains have mostly examined capability, grounding or confidence, and few score hallucination as a first-class outcome. CARES includes dermoscopy and fundus images within its trustworthiness suite \citep{xia2024cares}. A preprint on histopathology reports pathologist-graded errors, hallucinations and clinical impact for four general-purpose multimodal models \citep{lami2026histopathology}, and a second study reports that a general-purpose model produces confident descriptions when given blank histology images \citep{hasan2026seeing}. A study of vision-language models on chest radiography, brain MRI and dermatology shows that frontier models invent diagnoses when no image is supplied, and that the invented diagnosis depends on the demographic descriptors in the prompt \citep{vohra2026hearsay}. An audit of six large vision-language models for fundus images assesses dependence on text, logical consistency and the fidelity of chain-of-thought reasoning \citep{zhang2026fundus}. These studies establish that hallucination in non-radiological domains is real and measurable. They differ from one another in task, in grading and in the models examined, so that their results cannot be compared directly. A single protocol applied to several of these domains, with both domain-adapted and general-purpose models, is what the present benchmark provides.
% NOTE: lami2026histopathology was read through the medRxiv API abstract by an earlier session (lit_update.bib); the full text was not read. Cite the findings only at the level of the abstract.

\subsection{Severity-aware evaluation}
\label{subsec:rw_severity}

Two works grade errors by their clinical consequence, and they are the closest antecedents of the metrics proposed here. The MediHall Score assigns each statement to one of five levels, namely catastrophic, critical, attribute, prompt-induced and minor hallucination, with scores of 0.0, 0.2, 0.4, 0.6 and 0.8, and a score of 1.0 for a correct statement. Scores are averaged over the sentences of a report and over the answers of a set \citep{chen2024medhallmark}. The score is therefore normalised by all statements, correct ones included, and is sensitive to the frequency of errors. CREST defines twelve error types for chest radiograph report impressions, in four groups: incorrect affirmations or negations, mishandled uncertainty, omitted findings or negative evidence, and fabricated findings. Physicians assigned each type a high, medium or low clinical risk, which is mapped to weights of 3, 2 and 1 and normalised to the interval from 0 to 1. A companion measure reports the share of errors that are high risk \citep{crest2025medrxiv}.

The present work differs from these scores in three respects, and shares an important property with the first. Like the MediHall Score, \SMHR{} is normalised by all units and not by the errors alone; its difference lies in the unit, which is the atomic claim extracted from the response and not the sentence. Because both scores are weighted error rates over the same classified units, they are expected to order models similarly, and a direct comparison on the present data is reported in Section~\ref{sec:results}. First, the taxonomy is defined in terms of the relation between a claim and the image, includes omission as an explicit class, and is applied unchanged to four modalities. Second, \SRI{} expresses the \SMHR{} of a model relative to a fixed general-purpose reference in the same modality, which permits a comparison of specialist and general models across modalities with different base rates. Third, the class weights are assigned by the authors. This is a limitation relative to CREST, whose risk levels were assigned by physicians. A sensitivity analysis over the weights is reported in Section~\ref{sec:results}, and it can establish only the stability of the conclusions with respect to the weights, not their clinical validity.
% NOTE: Section results must contain the SMHR / MediHall / CREST comparison (claim_audit_report.md, C075-C078): SMHR and the MediHall Score ranked the 12 models identically by construction; SMHR and CREST rho = 0.64 (permutation p = 0.03).

\subsection{Domain adaptation and hallucination}
\label{subsec:rw_special}

Evidence on whether domain adaptation reduces hallucination is mixed. Wu et al.\ and Gu et al.\ report that medical models can be less robust than general models \citep{wu2024halluvqa,gu2024medvh}. EchoBench, which covers 20 modalities, reports that many medical-specific models exceed 95\% sycophancy \citep{yuan2025echobench}, and a benchmark of 5,000 items built from PathVQA, SLAKE and VQA-RAD found that medical fine-tuning tends to increase the tendency of a model to abandon a correct image-based answer under social pressure, with failure rates up to 75\% \citep{xu2025sycophancy}. A study of six models, three general-purpose and three medical, on three VQA datasets found that the models which hallucinated least were the most prone to yield to user pressure, and that none of the evaluated models of 7--8 billion parameters performed well on both grounding and pressure resistance \citep{aranya2026grounding}.

Paired designs, which compare a medical model with the general model from which it was derived, give a less uniform picture. A study of two such pairs, LLaVA-1.6 with LLaVA-Med and Gemma-3 with MedGemma, on four binary tasks in brain MRI, chest radiography, dermoscopy and histopathology found no reliable improvement from medical fine-tuning, with results that depended on the task and on the wording of the prompt \citep{mclaughlin2026knowledge}. A behavioural audit on brain MRI that included a base-versus-specialist pair found that medical adaptation improved the detection of tumours, while hallucination and abstention varied separately from accuracy \citep{sabbaghziarani2026confident}. These studies differ in task, in the way hallucination is measured and in the models included, and their results are not directly comparable with one another or with the present study. They indicate that a comparison of domain-adapted and general models should be reported as a measurement conditional on the models, the task and the metric, and that a design which does not contain base--specialist pairs cannot attribute a difference to the adaptation itself. The comparison in Section~\ref{sec:results} is formulated in this way: it compares each domain-adapted model with the best of nine general-purpose models and does not isolate the effect of medical adaptation.

\subsection{Calibration of stated confidence}
\label{subsec:rw_calibration}

Three of the four query types in the present protocol request a confidence value, which allows the calibration of the stated confidence to be examined alongside hallucination. Studies of medical VQA report that overconfidence persists across model families and is not removed by increasing scale or by prompting, that post-hoc calibration is more effective than prompt-based elicitation, and that signals from hallucination detection can improve calibration \citep{byun2026overconfidence}. A study of uncertainty methods over twelve VLMs and eight modalities reports that uncertainty is not a safety net for clinical VQA \citep{fazla2026uncertainty}, and a benchmark of calibration in medical VLMs examines classification across three modalities \citep{sharifdeen2026mvcbench}. The calibration experiments of the present study are reported in Section~\ref{subsec:res_calibration}. They vary the decoding setting, which, to our knowledge, was not varied in these studies.
% NOTE: Byun et al. was read at abstract level in this session. The statement about decoding settings follows claims_at_risk.md (section 5); re-run the literature search before submission.

\subsection{Mitigation of hallucination}
\label{subsec:rw_mitigation}

Mitigation strategies for VLMs act on the inputs, on the decoding procedure or on the model. Visual contrastive decoding is a training-free method that compares the output probabilities obtained from the original image with those obtained from a deliberately degraded version, and thereby steers generation towards content supported by the visual evidence \citep{leng2024vcd}. Retrieval-augmented generation supplies external text to the model at inference time \citep{lewis2020rag}. In the medical domain, MedHEval evaluates seven mitigation techniques that act on decoding or attention on 11 models, using radiological data \citep{chang2025medheval}, MedHallTune uses instruction tuning \citep{yan2025medhalltune}, and a single-call prompting method that removes social cues before the model answers has been proposed for sycophancy \citep{xu2025sycophancy}. A preprint on chest radiograph report generation reports that retrieval can itself induce hallucination when the retrieved context does not match the image \citep{idoko2026retrieval}. A mitigation experiment requires a baseline generated under the same decoding configuration as the mitigated arm, because a change in the decoding setting can alter the output independently of the strategy under test. The experiments in Section~\ref{sec:results} therefore report each strategy against two baselines, the benchmark response and a matched greedy response, and state the decoding configuration of every arm.

\subsection{Positioning of the present work}
\label{subsec:rw_position}

Severity-aware scoring, taxonomies of hallucination, benchmarks for medical hallucination, comparisons of domain-adapted with general models and studies of mitigation all have precedents, and none of them is presented here as new in itself. Localising the stage at which a hallucination arises is the subject of ClinHallu \citep{yang2026clinhallu} and is not attempted here. The contribution of the present work lies in the combination of a single taxonomy, a single query protocol and a single set of claim-level metrics applied to four imaging modalities for both domain-adapted and general-purpose models, in a reference-normalised index that supports comparison across modalities, in intervals that treat the image as the unit of independent evidence, and in a mitigation comparison in which each strategy is reported against two baselines. Among the works reviewed above, those that grade errors by clinical consequence were developed for report generation or for a single modality, and for several of the broader benchmarks the modality coverage and the grading scheme could not be established from the sources consulted (Table~\ref{tab:related}). A literature search should be repeated shortly before submission, since the number of relevant studies is increasing quickly.
